\section{An explicit general linear action and Lefschetz maps}
\label{sec:hook-action}

Let $\mu=(c,\tau)\vdash n$, put $r=|\tau|=n-c>0$, and set
\[
 N=\binom n2,\qquad S=\frac{r(r+1)}2,\qquad L=N-S+1.
\]
Write $K_{a,j}$ for the primitive harmonic-form model of the row hook
$(a,1^j)$ from Lemma~\ref{lem:hook-bridge}, and put
\[
 \mathcal M_\mu=\bigoplus_{a\ge1}
          \operatorname{Hom}_{S_n}(S^\mu,K_{a,r-1}).
\]
Its grading is the first row length $a$. We construct the representation
behind Theorem~\ref{thm:maximal-hook-series}, using the bounded-exponent
model of the ordinary coinvariant algebra.
The spaces $K_{a,r-1}$ are the canonical harmonic-form models of
stable row-hook multiplicity spaces. The auxiliary action is defined on $\mathcal M_\mu$.
In first-row degree $a$, the ordinary $\operatorname{GL}_k$-isotypic
component is $\mathbb S_{(a,1^{r-1})}K^k\otimes\mathcal M_\mu[a]$,
with the row group acting on the first factor.

Let $C_n=K[x_1,\ldots,x_n]/(K[x_1,\ldots,x_n]^{S_n}_+)$.
For $1\le d<n$, let $J_{n,d}\subset C_n$ be generated by the
determinants $\det(x_{i_a}^{\,b})_{1\le a,b\le d}$ over all
$d$-element label sets. Write $\Omega_n$ for the coinvariant algebra
of differential forms, apolar dual to $\SH_n$.
Proposition~\ref{prop:standalone-eulerian-foundations} proves the
bounded-exponent complement and its Foulkes character,
$\operatorname{Ann}_{C_n}(J_{n,d})=J_{n,n-d}$, and the
prefix-contraction image used below.

Let $U_c=K[t]_{<c}$ with ordered basis $1,t,\ldots,t^{c-1}$, and
identify $B_c=U_c^{\otimes n}$ with the bounded-exponent subspace
\[
 B_c=\operatorname{span}\{x_1^{a_1}\cdots x_n^{a_n}:0\le a_i<c\}
                  \subset K[x_1,\ldots,x_n].
\]
The group $\operatorname{GL}(U_c)$ acts on the tensor factors and
commutes with $S_n$. Transporting this action will give the action on $\mathcal M_\mu$.

\begin{lemma}[The bounded-exponent realization at the minimum descent number]
\label{lem:maximal-hook-bounded-model}
The quotient projection induces a graded $S_n$-equivariant isomorphism
\[
 (B_c)_{\mu'}\xrightarrow{\ \sim\ }
                  (C_n/J_{n,c})_{\mu'},
\]
with the quotient and determinant ideal just defined.
\end{lemma}
\begin{proof}
The straightening argument in Appendix~\ref{app:eulerian-foundations}
gives the surjection
$B_c\twoheadrightarrow C_n/J_{n,c}$. Applying the $\mu'$ central
idempotent preserves surjectivity. Schur--Weyl duality identifies the
graded $S^{\mu'}$ multiplicity in the source with
$\mathbb S_{\mu'}U_c$, whose character is
$s_{\mu'}(1,q,\ldots,q^{c-1})$.

Equations~\eqref{eq:standalone-eulerian-complement} and
\eqref{eq:standalone-descent-filtration} identify the target
multiplicity character as
\[
 \sum_{\substack{T\in\operatorname{SYT}(\mu')\\
                              \operatorname{des}(T)<c}}q^{\operatorname{maj}(T)}.
\]
The shape $\mu'$ has exactly $c$ rows, so each tableau has at least
$c-1$ descents. Principal specialization on $c$ variables identifies
the displayed sum, consisting precisely of tableaux with $c-1$
descents, with $s_{\mu'}(1,q,\ldots,q^{c-1})$.
Thus the surjection has equal source and target dimensions in every
degree, proving injectivity as well.
\end{proof}

Use the harmonic-form operators
\[
 \Theta_j=\sum_i\theta_i\partial_{x_i}^j,
 \qquad \iota=\sum_i x_i\partial_{\theta_i},
 \qquad d=\dd=\Theta_1.
\]
For a homogeneous $b\in(B_c)_{\mu'}$ of degree $h$, define
\begin{equation}\label{eq:explicit-maximal-hook-map}
 \Phi_\mu(b)=\frac{1}{N-h-S+r}\,
        \iota\Theta_1\Theta_2\cdots\Theta_r\,b(\partial)\Delta_n.
\end{equation}
The denominator is positive whenever the source component is nonzero.
Indeed,
$\mathbb S_{\mu'}U_c=\det(U_c)\otimes\mathbb S_{\tau'}U_c$,
so its largest polynomial degree is at most
$\binom c2+r(c-1)=N-S$. Consequently
$N-h-S\ge0$ and $N-h-S+r\ge r>0$.

\Needspace{17\baselineskip}
\begin{samepage}
\begin{theorem}[An explicit Schur-module realization]
\label{thm:maximal-hook-action}
The map \eqref{eq:explicit-maximal-hook-map} is a sign-intertwining
isomorphism
\[
 (B_c)_{\mu'}\xrightarrow{\ \sim\ }
                     \bigoplus_{a\ge1}(K_{a,r-1})_\mu,
 \qquad h\longmapsto a=L-h.
\]
On multiplicity spaces, it constructs an auxiliary
$\operatorname{GL}_c$ action for which
\[
                    \mathcal M_\mu\cong\mathbb S_{\tau'}U_c
\]
as representations, with the displayed grading given by the
principal cocharacter
$\rho(z)=\operatorname{diag}(1,z,\ldots,z^{c-1})$:
weight $a-1$ is the harmonic multiplicity space $\mathcal M_\mu[a]$.
Thus the graded character has the shift by $q$ in
Theorem~\ref{thm:maximal-hook-series}; the group acts across first-row degrees.
The action is obtained from the fixed bounded-exponent model and the
displayed differential map; no choice of a graded basis in the harmonic
multiplicity spaces is involved.
\end{theorem}
\end{samepage}
\begin{proof}
The perfect ordinary coinvariant pairing is
\[
 C_n[h]\otimes C_n[N-h]\longrightarrow K\Delta_n.
\]
The complementary annihilator identity
\eqref{eq:standalone-eulerian-annihilator} states
$\operatorname{Ann}_{C_n}(J_{n,r})=J_{n,c}$.
Consequently the pairing restricts to a perfect sign-valued pairing
between $C_n/J_{n,c}$ and $J_{n,r}$, reversing degree around $N$.

On $\Omega_n$, put $\delta_j=\sum_i x_i^j\partial_{\theta_i}$.
The contraction $\delta_r\cdots\delta_1:\Omega_n^r\to C_n$ has
image $J_{n,r}$, and its apolar adjoint is
$\Theta_1\cdots\Theta_r$. On label $\mu$, the largest possible
descent number is $r$. The Fields character~\cite[Theorem~5.2]{MRW} in exterior degree $r$
therefore has only the tableaux with exactly $r$ descents and has
polynomial degrees $\operatorname{maj}(T)-S$.
The character of $(J_{n,r})_\mu$ consists of those same tableaux,
with degrees $\operatorname{maj}(T)$. Thus the prefix contraction
is an isomorphism on that label, with degree shift $S$.
Its adjoint identifies the perfect pairing above with the concrete map
\[
 b\longmapsto \Theta_1\cdots\Theta_r b(\partial)\Delta_n
\]
from $(C_n/J_{n,c})_{\mu'}$ onto $(\mathrm{SH}_n^r)_\mu$, with
the sign twist. Its kernel on $C_n$ is exactly the annihilator
$J_{n,c}$, so the map factors through the stated quotient.

There is no label-$\mu$ component in exterior degree $r+1$.
Hence every form in $(\mathrm{SH}_n^r)_\mu$ is killed by $d$.
The Euler identity $d\iota+\iota d=(p+r)\operatorname{id}$ in
polynomial degree $p$ identifies that space with primitive forms of
exterior degree $r-1$, by $\iota/(p+r)$, with inverse $d$.
Taking $p=N-h-S$ proves the explicit formula and the isomorphism,
using the preceding bounded-exponent lemma.

Apolarity reverses polynomial degree. Reversing the monomial basis
converts this to principal weights, and a determinant twist removes
the full first column of $\mu'$. Explicitly, let
$J:U_c\to U_c$ send $t^j$ to $t^{c-1-j}$.
Transport to the target the following action on the source multiplicity
space $\mathbb S_{\mu'}U_c$:
\[
       g\longmapsto \det(g)^{-1}\mathbb S_{\mu'}(JgJ^{-1}).
\]
Since $\mu'=(1^c)+\tau'$, this is the polynomial irreducible
representation $\mathbb S_{\tau'}U_c$. If
$g=\operatorname{diag}(1,z,\ldots,z^{c-1})$, its eigenvalue on source
degree $h$ is
\[
 z^{n(c-1)-h-\binom c2}=z^{L-h-1}=z^{a-1}.
\]
Here $N-S=\binom c2+r(c-1)$.
This proves the grading shift by one. Normalizing the Specht pairing
or Vandermonde rescales the intertwiner and does not change the
transported action.
\end{proof}

\begin{corollary}[Lefschetz maps on harmonic multiplicities]
\label{cor:maximal-hook-lefschetz}
There are explicitly defined operators $E,F,H$ on $\mathcal M_\mu$
satisfying the $\mathfrak{sl}_2$ relations. They have grading changes
$+1,-1,0$, and
\[
 H|_{\mathcal M_\mu[a]}=(2(a-1)-r(c-1))\operatorname{id}.
\]
In particular, whenever $a\le1+r(c-1)/2$,
\[
 E^{\,2+r(c-1)-2a}:\mathcal M_\mu[a]
       \xrightarrow{\ \sim\ }
                    \mathcal M_\mu[2+r(c-1)-a].
\]
\end{corollary}
\begin{proof}
On $B_c$, use the usual tensor action of the irreducible
$\mathfrak{sl}_2$ module $U_c=\operatorname{Sym}^{c-1}K^2$:
\[
 E_B=\sum_i((c-1)x_i-x_i^2\partial_{x_i}),\qquad
 F_B=\sum_i\partial_{x_i},\qquad
 H_B=2\sum_i x_i\partial_{x_i}-n(c-1).
\]
These preserve the bounded-exponent space and commute with $S_n$.
Basis reversal exchanges $E_B$ and $F_B$ and negates $H_B$.
Thus the desired operators are
\[
 E=\Phi_\mu F_B\Phi_\mu^{-1},\qquad
 F=\Phi_\mu E_B\Phi_\mu^{-1},\qquad
 H=-\Phi_\mu H_B\Phi_\mu^{-1},
\]
on the sign-adjusted multiplicity spaces. The inverse means the unique
bounded-exponent preimage under the explicit isomorphism above.
The Cartan eigenvalue follows from $a=L-h$. Decomposition into finite
dimensional $\mathfrak{sl}_2$ modules gives the stated Lefschetz
isomorphisms. Since all row hooks are Vandermonde-generated, these
operators act on the same actual subspaces in both modules.
\end{proof}

